\documentclass[11pt]{article}
\usepackage[margin=0.78in]{geometry}
\usepackage[T1]{fontenc}
\usepackage[utf8]{inputenc}
\usepackage{amsmath,amssymb}
\usepackage{booktabs,longtable,array}
\usepackage[table]{xcolor}
\usepackage{enumitem}
\usepackage[hidelinks]{hyperref}

\newcommand{\wt}{\widetilde}
\newcommand{\Fcal}{\mathcal{F}}
\newcommand{\W}{\widetilde W}
\newcommand{\issue}[1]{\textbf{#1}}
\newcommand{\code}[1]{\texttt{#1}}
\newcolumntype{L}[1]{>{\raggedright\arraybackslash}p{#1}}

\definecolor{hard}{HTML}{B42318}
\definecolor{struct}{HTML}{875A00}
\definecolor{minor}{HTML}{3D5A80}

\title{Remaining-Points Check of the Updated Calculation}
\author{Check note}
\date{7 June 2026}

\begin{document}
\maketitle

\section*{Executive summary}

This note checks the full active body of the updated TeX/PDF, not only
the shortcut section.  The no-shortcut front part mostly passes direct finite
matrix checks: the preliminary identities, the no-shortcut first-passage
formula, the killed propagator, and the Chebyshev compact forms agree with
direct sums or absorbing-matrix resolvents to roundoff.  The front part still
has two local points worth adjusting: the compact Eq. (13) upper limit appears
to use the long-mode range if read literally after \(f_k,\alpha_k\) are
defined, and two identities on line 365 are off by \(1/2\).

The shortcut part has a separate downstream sign-propagation point.  The
updated manuscript has corrected the shortcut resolvent sign through Eq. (31),
but the correction has not been propagated consistently.  The first main
remaining shortcut point is Eq. (32): after summing the corrected Eq. (31), the
denominator should remain
\[
  1+z\beta(1-q)\wt W_u(u,z),
\]
not \(1-z\beta(1-q)\wt W_u(u,z)\).  Because Eq. (33), Eq. (40), Eq. (41),
and the final tail paragraph are downstream of Eq. (32), they inherit this
remaining sign mismatch.

The shortest way to see the point is:
\begin{align*}
\boxed{\text{correct Eq. (31)}} &\Longrightarrow
\boxed{\text{Eq. (32) denominator should be plus}}\\
&\Longrightarrow
\boxed{\text{Eq. (33) and Eq. (40) should have plus shortcut terms}}.
\end{align*}
For the symmetric antipodal case, the safest corrected downstream expression is
the closed form
\[
\boxed{
  \wt F_\rho(z)=
  \frac{aT_{L-\rho}(y)+U_{\rho-1}(y)+U_{L-\rho-1}(y)}
  {aT_L(y)+U_{L-1}(y)}
}
\]
where \(N=2L\), \(a=q/[\beta(1-q)]\), \(y=1+(1/z-1)/q\), and \(U_{-1}=0\).
The time-domain tail is controlled by the roots of
\[
  D(y)=aT_L(y)+U_{L-1}(y),
\]
not by a generic \(t\alpha_1^t\) term.

This check covers the active manuscript body up to \(\backslash\)end\{document\}
in the updated TeX.  Material after that point is commented-out scratch text and
is not treated as part of the active PDF.

\section*{Check scope}

The working instruction for this pass is:

\begin{quote}
Check the latest `Calculations.tex' and `Calculations.pdf' as a full active
manuscript, not only the shortcut section beginning around Eq. (29).  Check the
preliminary identities, no-shortcut first-passage formulae, killed propagator
formulae, compact modal notation, Chebyshev generating functions, and the
special antipodal self Green function.  For each formula, distinguish
numerically verified formulae from notationally fragile formulae, formulae that
differ if used literally, points that affect downstream conclusions, and
compile/self-containedness points.
\end{quote}

\section*{Pre-shortcut validation}

The following finite checks were run before evaluating the shortcut section.
They compare the printed formulae with direct finite sums, transition matrices,
absorbing submatrices, and resolvents.

\[
\begin{array}{lll}
\text{Check} & \text{Max difference} & \text{Interpretation}\\
\hline
\text{Eq. (1)--(4) preliminary identities} & 2.486900\times10^{-14} & \text{passes}\\
\text{Eq. (8) periodic propagator} & 5.551115\times10^{-16} & \text{passes}\\
\text{Eq. (9)--(12) first-passage time form} & 2.498002\times10^{-16} & \text{passes}\\
\text{Eq. (13), corrected compact upper }N/2 & 2.775558\times10^{-16} & \text{passes}\\
\text{Eq. (13), literal printed upper }N-1 & 4.811867\times10^{-1} & \text{differs}\\
\text{Eq. (16) killed propagator time form} & 1.642205\times10^{-15} & \text{passes}\\
\text{Eq. (19)--(20) killed propagator GF sums} & 7.993606\times10^{-15} & \text{passes}\\
\text{Eq. (21)--(22) Chebyshev modal sums} & 5.551115\times10^{-15} & \text{passes}\\
\text{Eq. (27)--(28) self Green functions} & 3.552714\times10^{-15} & \text{passes}
\end{array}
\]

Thus the front part is broadly consistent.  Its main formula-level point is the
compact-mode indexing in Eq. (13).  Once \(f_k\) and \(\alpha_k\) have been
defined for \(k=1,\ldots,N/2\), the compact generating-function sum must use
\(\sum_{k=1}^{N/2}\), not \(\sum_{k=1}^{N-1}\).  The Chebyshev expression on
the right-hand side agrees with the corrected compact sum.

\section*{Core corrected chain}

When the shortcut goes from \(u\) into the absorbing target \(v\), it is killing
inside the transient state space.  The corrected first-passage generating
function is
\begin{equation}
\boxed{
\wt F_{n_0}(v,z)=
\wt{\Fcal}_{n_0}(v,z)
+qz\,\wt W_{n_0}(u,z)
\left[1-\wt{\Fcal}_{u}(v,z)\right]\wt H(z)}
\label{eq:correct-F-H}
\end{equation}
with
\begin{equation}
\boxed{
  \wt H(z)=
  \frac{\beta(1-q)}{q}
  \frac{1}{1+z\beta(1-q)\wt W_u(u,z)}.}
\label{eq:correct-H}
\end{equation}
Thus the corrected version of Eq. (32) is
\begin{equation}
\boxed{
\wt F_{n_0}(v,z)=
\wt{\Fcal}_{n_0}(v,z)+
\frac{z\beta(1-q)\wt W_{n_0}(u,z)
\left[1-\wt{\Fcal}_{u}(v,z)\right]}
{1+z\beta(1-q)\wt W_u(u,z)}.}
\label{eq:correct-eq32}
\end{equation}

In time domain this gives, for \(t\geq 1\),
\begin{equation}
\boxed{
F_{n_0}(v,t)=\Fcal_{n_0}(v,t)
+q\sum_{t'=0}^{t-1}W_{n_0}(u,t-1-t')
\sum_{t''=0}^{t'}
\left[\delta_{t',t''}-\Fcal_u(v,t'-t'')\right]
H(t'').}
\label{eq:correct-eq40}
\end{equation}

\section*{Corrected direct-route form}

If one follows the original Eq. (41) route, \(H(0)\) must be split off from the
positive-time pole expansion:
\[
  h_0=H(0)=\frac{\beta(1-q)}q=\frac1a,
  \qquad H(c)=\sum_k c_k s_k^{c-1}\quad(c\geq1).
\]
The needed convolution kernels are
\[
  A^{(+)}_t(x;s)=\frac{s^{t-1}-x^{t-1}}{s-x},
  \qquad
  B^{(0)}_t(x,y)=\frac{x^{t-1}-y^{t-1}}{x-y},
\]
\[
  B^{(+)}_t(x,y;s)=
  [z^{t-3}]\frac{1}{(1-zx)(1-zy)(1-zs)}.
\]
Then the corrected long expression has four shortcut-convolution blocks:
\begin{align}
F_{n_0}(v,t)
&=
\sum_{\ell=1}^{L}f_\ell(n_0,v)\alpha_\ell^{t-1} \notag\\
&\quad
+qh_0
\left[
\sum_{r=1}^{N-1}G^{(1)}_r(n_0,u,v)\gamma_r^{t-1}
+
\sum_{r=1}^{L}g^{(2)}_r(n_0,u,v)\alpha_r^{t-1}
\right]\notag\\
&\quad
+q\sum_k c_k
\left[
\sum_{r=1}^{N-1}G^{(1)}_r(n_0,u,v)A^{(+)}_t(\gamma_r;s_k)
+
\sum_{r=1}^{L}g^{(2)}_r(n_0,u,v)A^{(+)}_t(\alpha_r;s_k)
\right]\notag\\
&\quad
-qh_0\sum_{\ell=1}^{L}f_\ell(u,v)
\left[
\sum_{r=1}^{N-1}G^{(1)}_r(n_0,u,v)B^{(0)}_t(\gamma_r,\alpha_\ell)
+
\sum_{r=1}^{L}g^{(2)}_r(n_0,u,v)B^{(0)}_t(\alpha_r,\alpha_\ell)
\right]\notag\\
&\quad
-q\sum_k c_k\sum_{\ell=1}^{L}f_\ell(u,v)
\left[
\sum_{r=1}^{N-1}G^{(1)}_r(n_0,u,v)B^{(+)}_t(\gamma_r,\alpha_\ell;s_k)
+
\sum_{r=1}^{L}g^{(2)}_r(n_0,u,v)B^{(+)}_t(\alpha_r,\alpha_\ell;s_k)
\right].
\label{eq:correct-hard-route}
\end{align}
This direct-route expression is useful for checking the original derivation, but
it should recompress to the closed form in the executive summary.  In that
closed form, for simple roots \(D(y_k)=0\), \(s_k=1-q+qy_k\),
\[
\boxed{
F_\rho(t)=\sum_k B_{\rho k}s_k^{t-1},
\qquad
B_{\rho k}=\frac{qN_\rho(y_k)}{D'(y_k)},}
\]
where \(N_\rho(y)=aT_{L-\rho}(y)+U_{\rho-1}(y)+U_{L-\rho-1}(y)\).

\section*{Checklist}

\small
\begin{longtable}{@{}L{0.08\linewidth}L{0.14\linewidth}L{0.24\linewidth}L{0.28\linewidth}L{0.10\linewidth}@{}}
\toprule
\textbf{ID} & \textbf{Location} & \textbf{Current point} &
\textbf{Suggested action} & \textbf{Type}\\
\midrule
\endfirsthead
\toprule
\textbf{ID} & \textbf{Location} & \textbf{Current point} &
\textbf{Suggested action} & \textbf{Type}\\
\midrule
\endhead

\issue{A0} & Eq. (13), line 251 &
The compact sum is printed as \(\sum_{k=1}^{N-1}\) even though \(f_k\) and
\(\alpha_k\) have just been defined with odd-mode compact indexing
\(k=1,\ldots,N/2\).  Used literally, this gives a different generating
function from the Chebyshev RHS/direct check. &
Use \(\sum_{k=1}^{N/2}\) in the compact form.  The Chebyshev expression on the
right agrees with this corrected compact sum. &
\textcolor{hard}{Formula-local}\\

\issue{A1} & Eq. (32), line 403 &
The denominator is still \(1-z\beta(1-q)\wt W_u(u,z)\). &
Use Eq. \eqref{eq:correct-eq32}; the denominator is
\(1+z\beta(1-q)\wt W_u(u,z)\). &
\textcolor{hard}{Main}\\

\issue{A2} & Eq. (33), line 408 &
The shortcut term is written with a minus sign, while the following definition
of \(H\) has the plus denominator. &
Use Eq. \eqref{eq:correct-F-H}: the term is
\(+qz\,\wt W_{n_0}(u,z)[1-\wt{\Fcal}_u(v,z)]\wt H(z)\). &
\textcolor{hard}{Main}\\

\issue{A3} & Eq. (40), line 458 &
The time-convolution inherits the minus sign from the Eq. (33) sign choice. &
Use Eq. \eqref{eq:correct-eq40}; the convolution term is positive. &
\textcolor{hard}{Main}\\

\issue{A4} & Eq. (41), lines 463--465 &
The long expression does not yet give the preferred corrected expression: the
baseline targets \(u\)
instead of \(v\), the global signs inherit Eq. (40), the \(g^{(2)}\) sector is
paired with \(\gamma_r\) in places where it should be paired with
\(\alpha_r\), and \(H(0)\) is not separated. &
Use the direct-route form \eqref{eq:correct-hard-route}, or preferably the
closed form and partial fractions above. &
\textcolor{hard}{Main}\\

\issue{A5} & Final paragraph, line 468 &
The conclusion that the tail is generically dominated by \(t\alpha_1^t\) is an
intermediate-mode conclusion, not a conclusion from the recompressed
generating function. &
The tail is \(B_{\rho1}s_1^{t-1}\) when the dominant root is simple; the
finite-time double-peak question is separate from this tail statement. &
\textcolor{hard}{Main}\\

\issue{A6} & Final paragraph, line 468 &
The proposed finite transition \(s_1=\gamma_1\) is not supported for finite
positive \(\beta\) in the corrected denominator. &
At \(s=\gamma_1\), the \(\phi(1/s)\) term is zero and the denominator is
\(q/[\beta(1-q)]>0\), so this crossing is not a root condition. &
\textcolor{hard}{Main}\\

\issue{B1} & Eq. (38)--(39), lines 428--451 &
The pole expansion of \(H\) is treated as if only positive-time poles matter.
The residue formula and root count are fragile. &
Keep \(H(0)=\beta(1-q)/q\) separate.  Write
\(c_k=qT_L(y_k)/D'(y_k)\) with \(D(y)=aT_L(y)+U_{L-1}(y)\), and solve
\(D(y)=0\) directly. &
\textcolor{struct}{Structural}\\

\issue{B2} & Figure caption, lines 436--438 &
The caption says \(q>1\) is required for a non-zero shortcut probability, while
the model uses \(0<q<1\); it also uses the pole figure in a way that could be
read as determining the finite-time peak transition. &
Change to \(0<q<1\), \(\beta>0\).  Treat the figure as an \(H\)-pole
visualization only, not as a standalone double-peak-transition criterion. &
\textcolor{struct}{Structural}\\

\issue{B3} & Line 441 and Eq. (39), line 448 &
The text says \(k=1,\ldots,N/2-1\) in the denominator, while the previous
caption says \(M=N/2\); the coefficient formula is much less robust than a
derivative-residue expression. &
Use \(D(y_k)=0\), \(s_k=1-q+qy_k\), and
\(c_k=qT_L(y_k)/D'(y_k)\).  Avoid claiming \(c_k=-\infty\) at an asymptote. &
\textcolor{struct}{Structural}\\

\issue{B4} & Line 365 &
Two of the three displayed trigonometric identities are off by \(1/2\). &
The correct identities are
\(\sum 1/(1-\cos(\pi k/N))=(N^2-1)/3\) and
\(\sum (-1)^k/(1-\cos(\pi k/N))=-(N^2+2)/6\). &
\textcolor{struct}{Structural}\\

\issue{B5} & Line 283 &
The text says \(Q_{n_0}(n,t)=0\) at the absorbing boundary and as
\(t\to\infty\).  \(Q\) was defined as the unrestricted periodic propagator, so
this statement does not match \(Q\) as defined. &
This statement should refer to \(W\), the killed propagator. &
\textcolor{struct}{Structural}\\

\issue{B6} & Eq. (25), line 288 &
The compact form has a dangling \(G^{(2)}_k\) term without a displayed sum,
while the later equivalent form uses the correct compact \(\alpha_k\) sum. &
Add the missing sum or replace this line with the later
\(\gamma_k/\alpha_k\) form. &
\textcolor{struct}{Structural}\\

\issue{B8} & Line 401 &
The text says Eq. (30) is summed, but the needed special case is the simplified
Eq. (31) with \(m=v\). &
Reference the simplified equation directly before deriving Eq. (32). &
\textcolor{struct}{Structural}\\

\issue{B9} & Line 257 &
The first displayed sentence calls \(\sum_{s=1}^t\Fcal_{n_0}(m,s)\) a survival
probability, but this quantity is the cumulative absorption probability up to
time \(t\). &
Call the first quantity cumulative absorption; reserve survival probability for
\(R_{n_0}(t)=1-\sum_{s=0}^t\Fcal_{n_0}(m,s)\). &
\textcolor{struct}{Structural}\\

\issue{C1} & Line 1; lines 65--68 &
The attachment is not self-contained: direct compilation needs local
style files and the figure.  Also \code{idelinks} appears to be a typo and
\code{customcolours} is loaded twice. &
Either include the style files and \code{phifunction.pdf}, or replace them with
portable definitions.  Move link options to \code{hyperref}. &
\textcolor{minor}{Compile}\\

\issue{C2} & Line 189 &
There is an empty label, \(\backslash\)label\{eq:\}. &
Replace it with a meaningful unique label or remove it. &
\textcolor{minor}{Notation}\\

\issue{C3} & Line 220 &
There is an extra closing parenthesis, and the condition says \(n_0\neq n\)
where the context is the absorbing target. &
Remove the extra parenthesis and clarify the condition as starting away from
the absorbing site, e.g. \(n_0\neq m\), as appropriate. &
\textcolor{minor}{Notation}\\

\issue{C4} & Lines 293--294 &
The arguments are written \(n_0,n,,m\). &
Remove the doubled comma. &
\textcolor{minor}{Typo}\\

\issue{C5} & Line 370 &
The notation switches to \(\wt{\mathcal W}\) while the rest of the manuscript
uses \(\wt W\). &
Use a consistent killed-propagator symbol unless a distinct object is intended. &
\textcolor{minor}{Notation}\\

\issue{C6} & Lines 438, 441, 451 &
Small text edits: \code{correspoding}, \code{dominator}, and
\code{an\$s\_k}. &
Correct to \code{corresponding}, \code{denominator}, and ``and
\(s_k\neq\alpha_k\)''. &
\textcolor{minor}{Typo}\\

\issue{C7} & Lines 171 and 433 &
Minor active-text copy points: ``preliminaries identities'' and the sentence
``the denominator ... has \(M\) real zeros ... and are shown graphically''. &
Use ``preliminary identities'' or ``preliminaries: identities'', and rewrite
the line 433 sentence so the subject agrees, e.g. ``the zeros are shown''. &
\textcolor{minor}{Typo}\\

\bottomrule
\end{longtable}
\normalsize

\section*{Numerical checks}

\subsection*{Shortcut sign benchmark}

The direct benchmark is the finite stochastic shortcut matrix: starting from
the lazy ring transition matrix, move
\(\lambda=\beta(1-q)\) from the \(u\to u\) self-loop to the directed shortcut
\(u\to v\).  Then compute the resolvent and the first-passage ratio
\[
  \wt F_{n_0}(v,z)=\frac{\wt S_{n_0,v}(z)}{\wt S_{v,v}(z)}.
\]
Over the checked grid, the maximum differences were:
\[
\begin{array}{ll}
\text{updated/original Eq. (32) sign vs stochastic matrix} & 3.669746\times10^{-1}\\
\text{corrected plus/plus Eq. (32) vs stochastic matrix} & 2.775558\times10^{-16}\\
\text{literal Eq. (29) ratio vs original Eq. (32)} & 1.318390\times10^{-16}\\
\text{stochastic-sign Eq. (29) ratio vs corrected Eq. (32)} & 5.551115\times10^{-17}\\
\text{Eq. (19) } W_{n_0}(u,z)\text{ vs absorbing matrix} & 1.043610\times10^{-14}\\
\text{Eq. (28) } W_u(u,z)\text{ vs absorbing matrix} & 1.687539\times10^{-14}
\end{array}
\]
This isolates the point to the shortcut sign propagation, not to the \(W\)
inputs.

\subsection*{Identity check}

For \(N=10\),
\[
  \sum_{k=1}^{N-1}\frac{1}{1-\cos(\pi k/N)}=33,
  \qquad
  \frac{2N^2+1}{6}=33.5,
\]
and
\[
  \sum_{k=1}^{N-1}\frac{(-1)^k}{1-\cos(\pi k/N)}=-17,
  \qquad
  \frac{1-N^2}{6}=-16.5.
\]
So the second and third identities displayed on line 365 are not exact.

\subsection*{Pre-shortcut check summary}

The pre-shortcut checks show that the main no-shortcut formulae are numerically
sound once the compact index in Eq. (13) is corrected:
\[
\begin{array}{ll}
\text{preliminary identities Eq. (1)--(4)} & 2.486900\times10^{-14}\\
\text{first-passage Eq. (9)--(13), corrected compact index} & 2.775558\times10^{-16}\\
\text{killed propagator Eq. (16), Eq. (19)--(20)} & 7.993606\times10^{-15}\\
\text{Chebyshev forms Eq. (21)--(22), Eq. (27)--(28)} & 5.551115\times10^{-15}
\end{array}
\]
The literal printed Eq. (13) compact upper limit \(N-1\), however, differs from
the finite-matrix benchmark by up to \(4.811867\times10^{-1}\) on the check
grid.

\section*{Shortest collaborator-facing statement}

A concise way to describe the situation is:

\begin{quote}
I checked the active manuscript body, including the formulae before the
shortcut section.  The no-shortcut first-passage and killed-propagator formulae
agree with finite-matrix checks to roundoff once the compact Eq. (13) index is
read as \(k=1,\ldots,N/2\).  The earlier local points I see are therefore the
printed compact upper limit in Eq. (13), the two line-365 identities, and some
notation/text points.

In the shortcut section, the updated resolvent sign is correct through the
simplified absorbing case, but the first-passage formula still carries the previous
denominator sign.  Once the denominator in Eq. (32) is corrected, Eq. (33) and
the time-convolution Eq. (40) also need the plus shortcut term.  The long
Eq. (41) should be treated with care as a tail conclusion in its present form:
after the
\(H(0)\) term is separated and the modes are recombined, the symmetric case
reduces to the closed form with denominator \(aT_L(y)+U_{L-1}(y)\), and the
tail is controlled by the dominant root of that denominator.
\end{quote}

\end{document}
